% GENERATED FILE. Edit canonical lessons, then run npm run generate:books.
% canonical-source-hash: fnv1a64:a1cca94a92903c10
% canonical-lessons: SA-C119-darvi, SA-C119-katahah, SA-C119-patram, SA-C119-kupi, SA-C119-chatrika, SA-R119-first-pass-nature-and-animals, SA-R119-second-pass-animals-the-body-and-the-kitchen

\chapter{Ladle, Frying pan, Vessel, Bottle, Small umbrella}
\label{ch:sa-ladle-frying-pan-vessel-bottle-small-umbrella}
\hlchaptermodality{119}

\begin{chapteropening}
\textbf{By the end of this chapter:} \emph{I can say a ladle, a frying pan, a vessel, a bottle and a small umbrella.}

Complete the last lesson of chapter 119: i can say a ladle, a frying pan, a vessel, a bottle and a small umbrella.
\end{chapteropening}

\section[\texorpdfstring{\sk{दर्वी} (darvī)}{दर्वी (darvī)}]{\sk{दर्वी} (darvī) — a ladle}

\label{lesson:SA-C119-darvi}



\begin{quote}
Before the new one: say the Sanskrit for a chair, then the Sanskrit for a board.
\end{quote}

\subsection*{You'll want to know: \sk{दर्वी}}
\textbf{\sk{दर्वी}} — \emph{darvī} — \textquotedblleft{}a ladle\textquotedblright{}.

\begin{grammarlens}[title={What you've built}]
One more word for this chapter.
\end{grammarlens}

\subsection*{Your turn}
\begin{itemize}
\raggedright
  \item \emph{Say it:} \emph{darvī}
  \item \emph{Say it:} \emph{darvī}, once more
  \item \emph{Say it:} say \emph{darvī}
  \item \emph{Recall:} say the Sanskrit for a chair, then the Sanskrit for a board, then say \emph{darvī} again
\end{itemize}

\begin{tcolorbox}[breakable,colback=teal!4,colframe=teal!35!black,title={Before you move on}]
What does \sk{दर्वी} mean? (\textquotedblleft{}A ladle\textquotedblright{}.) Say it once more.
\end{tcolorbox}
\section[\texorpdfstring{\sk{कटाहः} (kaṭāhaḥ)}{कटाहः (kaṭāhaḥ)}]{\sk{कटाहः} (kaṭāhaḥ) — a frying pan}

\label{lesson:SA-C119-katahah}



\begin{quote}
Before the new one: say the Sanskrit for a board, then the Sanskrit for a ladle.
\end{quote}

\subsection*{You'll want to know: \sk{कटाहः}}
\textbf{\sk{कटाहः}} — \emph{kaṭāhaḥ} — \textquotedblleft{}a frying pan\textquotedblright{}.

\begin{grammarlens}[title={What you've built}]
One more word for this chapter.
\end{grammarlens}

\subsection*{Your turn}
\begin{itemize}
\raggedright
  \item \emph{Say it:} \emph{kaṭāhaḥ}
  \item \emph{Say it:} \emph{kaṭāhaḥ}, once more
  \item \emph{Say it:} say \emph{kaṭāhaḥ}
  \item \emph{Recall:} say the Sanskrit for a board, then the Sanskrit for a ladle, then say \emph{kaṭāhaḥ} again
\end{itemize}

\begin{tcolorbox}[breakable,colback=teal!4,colframe=teal!35!black,title={Before you move on}]
What does \sk{कटाहः} mean? (\textquotedblleft{}A frying pan\textquotedblright{}.) Say it once more.
\end{tcolorbox}
\section[\texorpdfstring{\sk{पात्रम्} (pātram)}{पात्रम् (pātram)}]{\sk{पात्रम्} (pātram) — a vessel}

\label{lesson:SA-C119-patram}



\begin{quote}
Before the new one: say the Sanskrit for a ladle, then the Sanskrit for a frying pan.
\end{quote}

\subsection*{You'll want to know: \sk{पात्रम्}}
\textbf{\sk{पात्रम्}} — \emph{pātram} — \textquotedblleft{}a vessel\textquotedblright{}.

\begin{grammarlens}[title={What you've built}]
One more word for this chapter.
\end{grammarlens}

\subsection*{Your turn}
\begin{itemize}
\raggedright
  \item \emph{Say it:} \emph{pātram}
  \item \emph{Say it:} \emph{pātram}, once more
  \item \emph{Say it:} say \emph{pātram}
  \item \emph{Recall:} say the Sanskrit for a ladle, then the Sanskrit for a frying pan, then say \emph{pātram} again
\end{itemize}

\begin{tcolorbox}[breakable,colback=teal!4,colframe=teal!35!black,title={Before you move on}]
What does \sk{पात्रम्} mean? (\textquotedblleft{}A vessel\textquotedblright{}.) Say it once more.
\end{tcolorbox}
\section[\texorpdfstring{\sk{कूपी} (kūpī)}{कूपी (kūpī)}]{\sk{कूपी} (kūpī) — a bottle, a flask}

\label{lesson:SA-C119-kupi}



\begin{quote}
Before the new one: say the Sanskrit for a frying pan, then the Sanskrit for a vessel.
\end{quote}

\subsection*{You'll want to know: \sk{कूपी}}
\textbf{\sk{कूपी}} — \emph{kūpī} — \textquotedblleft{}a bottle, a flask\textquotedblright{}.

\begin{grammarlens}[title={What you've built}]
One more word for this chapter.
\end{grammarlens}

\subsection*{Your turn}
\begin{itemize}
\raggedright
  \item \emph{Say it:} \emph{kūpī}
  \item \emph{Say it:} \emph{kūpī}, once more
  \item \emph{Say it:} say \emph{kūpī}
  \item \emph{Recall:} say the Sanskrit for a frying pan, then the Sanskrit for a vessel, then say \emph{kūpī} again
\end{itemize}

\begin{tcolorbox}[breakable,colback=teal!4,colframe=teal!35!black,title={Before you move on}]
What does \sk{कूपी} mean? (\textquotedblleft{}A bottle, a flask\textquotedblright{}.) Say it once more.
\end{tcolorbox}
\section[\texorpdfstring{\sk{छत्रिका} (chatrikā)}{छत्रिका (chatrikā)}]{\sk{छत्रिका} (chatrikā) — a small umbrella}

\label{lesson:SA-C119-chatrika}



\begin{quote}
Before the new one: say the Sanskrit for a vessel, then the Sanskrit for a bottle.
\end{quote}

\subsection*{You'll want to know: \sk{छत्रिका}}
\textbf{\sk{छत्रिका}} — \emph{chatrikā} — \textquotedblleft{}a small umbrella\textquotedblright{}.

\begin{grammarlens}[title={What you've built}]
One more word for this chapter.
\end{grammarlens}

\subsection*{Your turn}
\begin{itemize}
\raggedright
  \item \emph{Say it:} \emph{chatrikā}
  \item \emph{Say it:} \emph{chatrikā}, once more
  \item \emph{Say it:} say \emph{chatrikā}
  \item \emph{Recall:} say the Sanskrit for a vessel, then the Sanskrit for a bottle, then say \emph{chatrikā} again
\end{itemize}

\begin{tcolorbox}[breakable,colback=teal!4,colframe=teal!35!black,title={Before you move on}]
What does \sk{छत्रिका} mean? (\textquotedblleft{}A small umbrella\textquotedblright{}.) Say it once more.
\end{tcolorbox}
\section[(dialogue)]{Nature and animals — a first pass}

\label{lesson:SA-R119-first-pass-nature-and-animals}



\begin{quote}
Say the words for a bottle and a small umbrella.
\end{quote}

\subsection*{Your turn}
Say these aloud:

\begin{itemize}
\raggedright
  \item a mountain, a waterfall, a wave, lightning, thunder
  \item a storm, mist, a branch, a leaf, a flower
  \item a lotus, a star, the sun, a bear, a jackal
  \item a rabbit, a donkey, a boar, a bull, a buffalo
  \item a sheep, a pigeon, an owl, a heron, a bee
\end{itemize}

\begin{tcolorbox}[breakable,colback=teal!4,colframe=teal!35!black,title={Before you move on}]
A mountain? (\textbf{\sk{गिरिः}}, \emph{giriḥ}.) A waterfall? (\textbf{\sk{प्रपातः}}, \emph{prapātaḥ}.) A wave? (\textbf{\sk{तरंगः}}, \emph{taraṅgaḥ}.) Lightning? (\textbf{\sk{विद्युत्}}, \emph{vidyut}.) Thunder? (\textbf{\sk{गर्जनम्}}, \emph{garjanam}.) A storm? (\textbf{\sk{वात्या}}, \emph{vātyā}.) Mist? (\textbf{\sk{नीहारः}}, \emph{nīhāraḥ}.) A branch? (\textbf{\sk{शाखा}}, \emph{śākhā}.) A leaf? (\textbf{\sk{पर्णम्}}, \emph{parṇam}.) A flower? (\textbf{\sk{कुसुमम्}}, \emph{kusumam}.) A lotus? (\textbf{\sk{पद्मम्}}, \emph{padmam}.) A star? (\textbf{\sk{नक्षत्रम्}}, \emph{nakṣatram}.) The sun? (\textbf{\sk{सूर्यः}}, \emph{sūryaḥ}.) A bear? (\textbf{\sk{भल्लूकः}}, \emph{bhallūkaḥ}.) A jackal? (\textbf{\sk{शृगालः}}, \emph{śṛgālaḥ}.) A rabbit? (\textbf{\sk{शशकः}}, \emph{śaśakaḥ}.) A donkey? (\textbf{\sk{गर्दभः}}, \emph{gardabhaḥ}.) A boar? (\textbf{\sk{वराहः}}, \emph{varāhaḥ}.) A bull? (\textbf{\sk{वृषभः}}, \emph{vṛṣabhaḥ}.) A buffalo? (\textbf{\sk{महिषः}}, \emph{mahiṣaḥ}.) A sheep? (\textbf{\sk{मेषः}}, \emph{meṣaḥ}.) A pigeon? (\textbf{\sk{कपोतः}}, \emph{kapotaḥ}.) An owl? (\textbf{\sk{उलूकः}}, \emph{ulūkaḥ}.) A heron? (\textbf{\sk{बकः}}, \emph{bakaḥ}.) A bee? (\textbf{\sk{भ्रमरः}}, \emph{bhramaraḥ}.)
\end{tcolorbox}
\section[(dialogue)]{Animals, the body and the kitchen — a second pass}

\label{lesson:SA-R119-second-pass-animals-the-body-and-the-kitchen}



\begin{quote}
Say the words for a fly and a mosquito.
\end{quote}

\subsection*{Your turn}
Say these aloud:

\begin{itemize}
\raggedright
  \item a fly, a mosquito, an ant, a butterfly, a rooster
  \item a crane, a cheek, the chin, an elbow, a wrist
  \item a thumb, the chest, the waist, an eyebrow, the skin
  \item an ankle, a finger, a room, a chair, a board
  \item a ladle, a frying pan, a vessel, a bottle, a small umbrella
\end{itemize}

\begin{tcolorbox}[breakable,colback=teal!4,colframe=teal!35!black,title={Before you move on}]
A fly? (\textbf{\sk{मक्षिका}}, \emph{makṣikā}.) A mosquito? (\textbf{\sk{मशकः}}, \emph{maśakaḥ}.) An ant? (\textbf{\sk{पिपीलिका}}, \emph{pipīlikā}.) A butterfly? (\textbf{\sk{चित्रपतंगः}}, \emph{citrapataṅgaḥ}.) A rooster? (\textbf{\sk{कुक्कुटः}}, \emph{kukkuṭaḥ}.) A crane? (\textbf{\sk{सारसः}}, \emph{sārasaḥ}.) A cheek? (\textbf{\sk{कपोलः}}, \emph{kapolaḥ}.) The chin? (\textbf{\sk{चिबुकम्}}, \emph{cibukam}.) An elbow? (\textbf{\sk{कूर्परः}}, \emph{kūrparaḥ}.) A wrist? (\textbf{\sk{मणिबन्धः}}, \emph{maṇibandhaḥ}.) A thumb? (\textbf{\sk{अंगुष्ठः}}, \emph{aṅguṣṭhaḥ}.) The chest? (\textbf{\sk{वक्षः}}, \emph{vakṣaḥ}.) The waist? (\textbf{\sk{कटिः}}, \emph{kaṭiḥ}.) An eyebrow? (\textbf{\sk{भ्रूः}}, \emph{bhrūḥ}.) The skin? (\textbf{\sk{त्वक्}}, \emph{tvak}.) An ankle? (\textbf{\sk{गुल्फः}}, \emph{gulphaḥ}.) A finger? (\textbf{\sk{अंगुलिः}}, \emph{aṅguliḥ}.) A room? (\textbf{\sk{कक्षः}}, \emph{kakṣaḥ}.) A chair? (\textbf{\sk{आसन्दी}}, \emph{āsandī}.) A board? (\textbf{\sk{फलकम्}}, \emph{phalakam}.) A ladle? (\textbf{\sk{दर्वी}}, \emph{darvī}.) A frying pan? (\textbf{\sk{कटाहः}}, \emph{kaṭāhaḥ}.) A vessel? (\textbf{\sk{पात्रम्}}, \emph{pātram}.) A bottle? (\textbf{\sk{कूपी}}, \emph{kūpī}.) A small umbrella? (\textbf{\sk{छत्रिका}}, \emph{chatrikā}.)
\end{tcolorbox}
